\subsection{Relation to earlier work}
\label{sec:literature}

The standard morphisms and palindromic extensions used here belong to
classical Sturmian combinatorics. Justin's functional identity relates
iterated palindromic closure to composition of standard
morphisms~\cite[Lemma~2.1]{Justin2005}. Its left and right forms, with the
same conventions used below, appear in Perrin and
Reutenauer~\cite[Section~2]{PerrinReutenauer2023}. The incidence formulas
of Berth\'e, de Luca and Reutenauer describe the completed letter counts
of central words and their standard factorizations~\cite[Section~3]{BertheDeLucaReutenauer2008}.
These results supply the word transport and the primitive-count test in
our application. The involution obtained by reversing a directive in
\cite{BertheDeLucaReutenauer2008} changes the central word by Christoffel
duality; our reflection instead acts on a positive matrix orbit while the
directive remains fixed.

The language that appears in the answer is also familiar. If
$\theta(a)=ab$ and $\theta(b)=ba$, then
$(ab\mid ba)^*=\theta(\{a,b\}^*)$ is the image of the Thue--Morse
substitution. Reutenauer and Vuillon use this substitution together with
palindromic and antipalindromic closure to construct Markov
numbers~\cite{ReutenauerVuillon2017}: in their coding a Markov number is two more than the
length of $\operatorname{Pal}(\theta(\operatorname{Pal}(av)))$, and the injectivity of this
length is equivalent to the uniqueness conjecture~\cite[Theorem~1]{ReutenauerVuillon2017}.
\Cref{prop:thue-morse} identifies that length plus two, the length of the Christoffel word
$a\operatorname{Pal}(\theta(F_a(z)))b$, with our Gram label, for every palindrome $z$.
Abram, Lapointe and Reutenauer
extend the construction to whole Markov triples~\cite{AbramLapointeReutenauer2020}.
Our conclusion concerns the role of this language in a different
problem: it exhausts the finite words whose positive terminal vectors
are carried to another such vector by the prescribed rational reflection.
The proof establishes this exhaustion uniformly for the reflection of every
palindrome, before imposing palindromicity on the reflected words.

Palindromic symmetry is well established in the geometry of two-generator
groups. Kassel and Reutenauer prove that a cyclically reduced basis with
both words of odd length has a unique palindromic conjugate
basis~\cite{KasselReutenauer2007}. Gilman and Keen relate palindromic
words to axes perpendicular to the common perpendicular of the generator
axes~\cite{GilmanKeen2009}. These results concern primitive bases or
geometric axes. The decoder here retains the positive word, the finite
terminal vector and the fixed rational reflection, and applies to every
finite word, including nonprimitive ones.

The distinction between a full matrix and one of its entries matters
when connecting the result to Markov arithmetic. Reutenauer's
Christoffel correspondence is a bijection with whole proper Markov
triples~\cite[Theorem~1]{Reutenauer2009}. More recently, Veselov identifies
Markov fractions with Cohn matrix indices and describes continued
fractions through a mirrored Conway topograph~\cite[Sections~3--4]{Veselov2026}.
Jouteur, Paris-Romaskevich and Thomas characterize modular involutions
exchanging rational pairs and study Springborn operations, including a
$q$-deformed midpoint formula for Farey neighbors~\cite[Sections~5--7]{JouteurParisRomaskevichThomas2026}.
Those constructions retain a tree position, an arithmetic regularity
condition or a whole rational fraction. They do not identify which
arbitrary positive words survive reflection in the orbit considered here.

There are also injectivity theorems for restricted word languages.
Lapointe and Reutenauer prove that a Cohn matrix entry is strictly
increasing in military order on the language of one fixed Sturmian
sequence~\cite[Corollary~1]{LapointeReutenauer2021}. Our comparison allows
arbitrary palindromes, whose completions need not lie in one such
language. Their Gram labels are positive integers, and are Markov
numbers under the additional Christoffel hypothesis. Equal labels and a
specified root midpoint recover a full Gram congruence, which brings the
arbitrary-word decoder into play. This yields an exact classification of
the reflected pairs and shows why every nonempty pair fails the primitive
completed-count condition. No claim is made that an arbitrary equal-label
pair has the prescribed midpoint.
